% from ega0-21-fr.tex
\oldpage[0\textsubscript{IV-1}]{182}
est une extension finie de $K_0$. Pour tout idéal premier $\mathfrak p$ de
$A$, soit $k(\mathfrak p)$ le corps résiduel de $A$ en $\mathfrak p$. Alors,
pour tout corps $k$ tel que $k_0\subset k\subset K_0$ et tel que
$[K_0:k]<+\infty$, on a
\begin{equation}
\tag{21.9.8.1}
\label{0.21.9.8.1-fr}
\begin{split}
  &\operatorname{rg}_{k(\mathfrak p)}
    \bigl(\widehat\Omega^1_{(A/\mathfrak p)/k}
      \otimes_{A/\mathfrak p}k(\mathfrak p)\bigr)
    -\operatorname{rg}_{K_0}\Omega^1_{K_0/k}\\
  &\hspace{35mm}\geqslant
    \dim(A/\mathfrak p)
    +\operatorname{rg}_{k(\mathfrak p)}\Upsilon_{k(\mathfrak p)/k_0}.
\end{split}
\end{equation}
Supposons en outre que $[k_0:k_0^p]<+\infty$ (où $p$ est l'exposant
caractéristique de $k_0$). Alors on peut trouver un corps $k$ tel que
$k_0(K_0^p)\subset k\subset K_0$ et $[K_0:k]<+\infty$, pour lequel les deux
membres de (21.9.8.1) sont égaux.
\end{proposition}

\begin{proof}
Comme $A/\mathfrak p$ a même corps résiduel $K$ que $A$ et est complet, on
peut se borner au cas où $A$ est intègre et $\mathfrak p=0$~; on posera
alors $E=k(\mathfrak p)$, corps des fractions de $A$. Comme $K_0$ est
formellement lisse sur $k_0$ (19.6.1), il existe un $k_0$-monomorphisme
$K_0\to A$ faisant de $A$ une $K_0$-algèbre~; on sait en outre qu'il existe
une sous-$K_0$-algèbre $A_0$ de $A$, $K_0$-isomorphe à une algèbre de séries
formelles $K_0[[T_1,\ldots,T_r]]$ et telle que $A$ soit une $A_0$-algèbre
\emph{finie} (19.8.9), ce qui entraîne que $E$ est une extension finie du
corps des fractions $L_0=K_0((T_1,\ldots,T_r))$ de $A_0$. On a alors
$\Upsilon_{K_0/k_0}=0$ puisque $K_0$ est séparable sur $k_0$, et
$\deg.\operatorname{tr}_{L_0}E=0$~; en outre, comme $A$ est une
$A_0$-algèbre finie on a $\dim(A)=\dim(A_0)$ (16.1.5).

Si $p=1$, on a
$\operatorname{rg}_E(\widehat\Omega^1_{A/k}\otimes_AE)=r$ par (21.9.5) et
(17.1.4, (iii)), et $\Omega^1_{K_0/k}=0$ puisque $K_0$ est extension finie
séparable de $k$ (21.7.4)~; d'autre part $\Upsilon_{E/k_0}=0$, $E$ étant
extension séparable de $k_0$~; les deux membres de (21.9.8.1) sont donc égaux
dans ce cas.

Supposons maintenant que $p>1$~; si $B=k[[T_1,\ldots,T_r]]$, $A$ est une
$B$-algèbre finie, donc, en posant $B_1=k[[T_1^p,\ldots,T_r^p]]$,
$\widehat\Omega^1_{A/k}$ s'identifie à $\Omega^1_{A/B_1}$ (21.9.4), et en
désignant par $M$ le corps des fractions $k((T_1^p,\ldots,T_r^p))$, il
résulte de (20.5.9) que $\widehat\Omega^1_{A/k}\otimes_AE$ s'identifie à
$\Omega^1_{E/M}$~; l'inégalité (21.9.8.1) n'est autre alors que (21.9.6.2),
et la dernière assertion du corollaire résulte de (21.9.6, (ii)).
\end{proof}

% (continued in ega0-22-fr.tex, before its first page marker)
\section{Critères différentiels de lissité formelle et de régularité}
\label{section:0.22-fr}

Les principaux résultats de ce paragraphe sont~:
\begin{enumerate}
  \item[a)] Les critères (22.2.2) et (22.5.8), qui donnent des conditions
  nécessaires et suffisantes pour qu'un anneau local noethérien complet $A$
  contenant un corps $k$ soit formellement lisse sur $k$~; le critère
  (22.5.8) s'énonce sans faire intervenir de notions différentielles, et
  complète (19.6.6)~; l'énoncé (22.2.2), de nature un peu plus technique,
  permet de donner une \emph{classification}, pour $k$ fixé, des anneaux
  locaux $A$ considérés (22.2.5)~: ils sont donnés par leur dimension $n$ et
  leur corps résiduel $K$, sous la seule condition
  $n\geqslant\operatorname{rg}_K\Upsilon_{K/k}$.

  \item[b)] Le théorème (22.3.3), et le corollaire (22.7.6) du critère
  jacobien de Nagata (22.7.3), qui donnent des conditions permettant
  d'affirmer qu'un anneau \emph{localisé} d'un anneau local noethérien complet
  $A$ contenant un corps $k$ est géométriquement régulier sur $k$~; ils
  joueront un rôle important dans l'étude des propriétés «~fines~» des
  anneaux locaux faite au chap.~IV, §~7. On ne possède pour le moment aucune
  démonstration de (22.7.6) indépendante du critère jacobien de Nagata.

  \item[c)] Le critère jacobien de Zariski (22.6.7) et ses variantes, qui
  sont des conséquences faciles de (20.7.8).
\end{enumerate}

